% ===================================================================
%  કોષ્ટક નં. ૪ — પુખ્ત વય અને વૃદ્ધાવસ્થા.
%
%  Traduction, faite en 2026, en gujarati d'aujourd'hui, sur l'ido de
%  texto/io. Regles de la colonne : voir 10-kostak-01.tex. Deux
%  scenes : t04-c1-* la noce, t04-c2-* l'anniversaire du grand-pere.
%  Ecrit parigi.py ouvert, et chaque paire annoncee rendue par une
%  relative en « જે », comme au tableau 3.
%
%  UNE HESITATION LAISSEE DANS LE TEXTE, ET C'EST renvoji.py QUI L'A
%  VUE. Au bloc t04-c1-06-1, l'ido dit « la patrulo (10) di la
%  mariajitino », le pere de la MARIEE ; l'alinea avait ete ecrit
%  d'abord avec la mere du marie, puis corrige EN COURS DE PHRASE
%  sans effacer le premier essai — « વરની માતાના (10) પિતાએ — ના,
%  વહુના પિતાએ (10) ». Deux renvois (10) au lieu d'un. Aucune regle
%  de kolonoj.py ne pouvait le voir : la ligne etait du gujarati
%  correct, aucune lettre etrangere, aucune longueur en trop. C'est
%  le COMPTE DES RENVOIS qui a denonce le brouillon, exactement comme
%  les mots « Wait » et « Ordre » laisses dans les colonnes ourdoue
%  et tamoule avaient ete denonces autrement. Chaque outil de la
%  famille attrape une espece de negligence, et aucun n'attrape celle
%  des autres.
%
%  L'AUTOMNE DE CE TABLEAU EST « શરદ », ET C'EST LA QUATRIEME DES SIX
%  SAISONS. Le tableau 3 avait pose que le gujarati en compte six ;
%  celui-ci le verifie sur la seule des six que l'Europe reconnaisse
%  aussi. Mais l'hiver de la seconde scene oblige a choisir entre
%  હેમંત et શિશિર, que l'Europe confond : હેમંત est le froid sec de
%  decembre-janvier, શિશિર le froid de la fin, fevrier-mars, quand
%  les feuilles tombent. La planche grave decembre, de la neige sur
%  les toits et un calendrier ouvert au premier decembre : on ecrit
%  હેમંત. LE FRANCAIS A UN MOT POUR CE QUE LE GUJARATI PARTAGE EN
%  DEUX, et c'est la premiere fois du releve qu'une colonne ait a
%  couper une saison europeenne en son milieu.
%
%  ET LA NEIGE (19) EST « બરફ », QUI EST AUSSI LA GLACE — mot persan,
%  barf. Le Gujarat n'en voit pas, mais il en parle depuis le
%  Sultanat : les glaciers de l'Himalaya alimentaient les glacieres
%  des princes, et le mot est arrive avec la chose, portee a dos
%  d'homme. La colonne haoussa avait du composer « dusar ƙanƙara »,
%  la poudre de glace, faute d'avoir jamais vu tomber la neige. Ici
%  le mot etait la, mais il ne distingue pas ce qui tombe de ce qui
%  gele — et la planche demande les deux.
%
%  LE FILET DE LA PARENTE EST LE PLUS SERRE DU RELEVE, ET IL OBLIGE A
%  CHOISIR PLUS SOUVENT QUE LE PERSAN. Le persan comptait quatre
%  oncles ; le gujarati en compte autant et distingue en plus les
%  grands-parents par le cote :
%
%    દાદા / દાદી    le pere et la mere du PERE
%    નાના / નાની    le pere et la mere de la MERE
%    કાકા / કાકી    le frere du pere, et sa femme
%    મામા / મામી    le frere de la mere, et sa femme
%    ફોઈ / ફુવા     la soeur du pere, et son mari
%    માસી / માસા    la soeur de la mere, et son mari
%
%  La planche ne dit jamais de quel cote. La colonne a pris le cote
%  du PERE partout — દાદા le grand-pere, દાદી la grand-mere, કાકા
%  l'oncle —, pour la meme raison que la persane avait pris عمو : la
%  seconde scene est l'anniversaire du grand-pere PATERNEL, puisque
%  c'est chez les parents de Ioannes qu'on se rend. Une langue qui
%  distingue oblige a savoir ce que le texte d'origine n'a pas eu
%  besoin de dire — et celle-ci oblige deux fois plus souvent.
%
%  LE MEDECIN EST « ડૉક્ટર » ET LE REMEDE EST « દવા », ET C'EST LA
%  TROISIEME FOIS QUE LE RELEVE VOIT CETTE LIGNE PASSER AU MEME
%  ENDROIT. Le persan partageait sa medecine entre la peau et les
%  organes ; le haoussa entre l'ancienne medecine et la nouvelle —
%  magani le remede contre likita le medecin. Le gujarati fait le
%  meme partage que le haoussa, avec d'autres mots : દવા le remede,
%  વૈદ le medecin ayurvedique, ઓસડ le simple, તાવ la fievre, શરદી le
%  rhume, ઉધરસ la toux — tous d'ici ; ડૉક્ટર et તેની ચિઠ્ઠી
%  l'ordonnance (27) sont venus avec l'hopital. TROIS COLONNES, TROIS
%  LEXIQUES, UNE SEULE LOI : ce qui s'est ECRIT ailleurs arrive avec
%  ses mots.
%
%  ET LES ANANAS DU DESSERT (45) SONT « અનનાસ », QUI EST DU
%  PORTUGAIS — ananás, pris au tupi d'Amerique et debarque a Goa.
%  Avec મેજ la table et પાદરી le pretre, cela fait trois mots
%  portugais en quatre tableaux, et le troisieme est le plus long
%  voyage : de l'Amazonie a Lisbonne, de Lisbonne a Goa, de Goa a
%  cette table de noce.
% ===================================================================

%%K t04-apar-1 apar
\VUsaut{4.0mm}
\VUtitre{15.0pt}{60}{કોષ્ટક નં. ૪}
\VUsaut{1.6mm}
\VUtitre{11.4pt}{0}{પુખ્ત વય અને વૃદ્ધાવસ્થા.}
\VUsaut{3.2mm}

%%K t04-tit-1 sub
\VUcentre{11.4pt}{0}{\textit{પ્રથમ દૃશ્ય.}}
\VUsaut{1.6mm}
\VUtitre{11.4pt}{0}{લગ્નનું ભોજન.}
\VUsaut{2.6mm}

%%K t04-tit-2 sub
\VUpk{10.2pt}{40}{શરદ.}
\VUsaut{3.0mm}

%%K t04-c1-01-1 p
1. --- જુવાનો મોટા થયા છે. તેમાંનો એક, યોહાનેસ, લગ્ન કરે છે. અમે
\VUgras{લગ્નના ભોજનમાં} હાજર છીએ, જે એક જ મેજની આસપાસ વરકન્યાનાં બે
કુટુંબોને ભેગાં કરે છે.

%%K t04-c1-02-1 p
2. --- અમે \VUgras{શરદમાં} છીએ, \VUgras{સપ્ટેમ્બર મહિનામાં.}
\VUgras{ઓક્ટોબર} અને \VUgras{નવેમ્બર}નો ઠંડો પવન હજી ગામડાને તેનાં
લીલાં પાંદડાંથી વંચિત કરી શક્યો નથી. સાંજની હવા હૂંફાળી અને સુગંધી છે ;
તેથી રાતનું ભોજન બગીચા પાસેની \VUgras{ઓસરીમાં} \textsuperscript{(8)}
થાય છે.

%%K t04-c1-03-1 p
3. --- દૂર, \VUgras{ટેકરી} \textsuperscript{(3)} પર, યોહાનેસનાં
માતાપિતાનું ઘર છે, હરિયાળીમાં છુપાયેલો સુઘડ \VUgras{મહેલ}
\textsuperscript{(1)} ; ઉત્તમ દ્રાક્ષાસવ આપતી સુંદર દ્રાક્ષવાડીઓ
ટેકરીના ઢોળાવને ઢાંકે છે. દ્રાક્ષખેડૂત તેમને ઉછેરે છે અને સંભાળે છે.

%%K t04-c1-04-1 p
4. --- દ્રાક્ષવાડીઓની ઉપર \VUgras{તેતરનું} \textsuperscript{(2)} એક
ટોળું પસાર થાય છે, જે \VUgras{શિકારરક્ષકના} \textsuperscript{(5)}
નજીક આવવાથી ડરી ગયું છે. તે, બગલમાં \VUgras{બંદૂક} અને ખભે
\VUgras{શિકારથેલી} લઈને, \VUgras{ઝાડીની} \textsuperscript{(4)} તપાસ કરે
છે પોતાના \VUgras{કૂતરા} \textsuperscript{(6)} સાથે. તે જમીનની દેખરેખ
રાખે છે અને ચોરશિકારીઓને શિકાર નાશ કરતાં અટકાવે છે.

%%K t04-c1-05-1 p
5. --- \VUgras{ઓસરીની} બહાર \VUgras{દ્રાક્ષવેલો} ચઢે છે, જેના પરથી
ઘાટા \VUgras{દ્રાક્ષના ઝૂમખાં} \textsuperscript{(7)} લટકે છે.

%%K t04-c1-06-1 p
6. --- \VUgras{મેજને} \textsuperscript{(16)} શણગારતાં સુંદર શરદનાં ફૂલ
\VUgras{ગુલદાઉદી} \textsuperscript{(9)} છે ; વહુના
\VUgras{પિતાએ} \textsuperscript{(10)} તેમાંનું એક પોતાના કોટની
બટનકાજમાં ખોસ્યું છે.

%%K t04-c1-07-1 p
7. --- ભોજન પૂરું થવા આવે છે. \VUgras{પીરસનારો}
\textsuperscript{(11)} મીઠાઈની થાળીઓ લાવવા લાગ્યો છે અને થોડી વારમાં
ભોજનનાં જુદાં જુદાં સાધનો ઉપાડી લેશે, \VUgras{ચમચા}
\textsuperscript{(14)}, \VUgras{કાંટા} \textsuperscript{(12)},
\VUgras{છરીઓ} \textsuperscript{(13)}, \VUgras{મોટી થાળીઓ}
\textsuperscript{(15)}, જેમની હવે જરૂર નથી.

%%K t04-tit-3 sub
\VUpk{10.2pt}{40}{સગાંવહાલાં.}
\VUsaut{3.0mm}

%%K t04-c1-08-1 p
8. --- બધાં આનંદમાં દેખાય છે, અને વરની \VUgras{માતા}
\textsuperscript{(17)} જે સ્મિત સાથે પોતાનાં સંતાનોને જુએ છે તે તેમનું
સુખ સાબિત કરે છે.

%%K t04-c1-09-1 p
9. --- આ લગ્ન દેશનાં બે ધનવાન કુટુંબોને જોડે છે. યોહાનેસ,
\VUgras{વર} \textsuperscript{(18)}, મોટા જમીનદારનો \VUgras{દીકરો} છે
અને હવે તે કુશળ કારખાનેદારનો \VUgras{જમાઈ} બને છે.

%%K t04-c1-10-1 p
10. --- \VUgras{વરકન્યા} જુવાન છે, છતાં તેમની \VUgras{સગાઈ} ઘણા વખતથી
થઈ ગઈ હતી. \VUgras{જુવાન સ્ત્રી} \textsuperscript{(19)}, માર્થા,
નારંગીનાં ફૂલથી શણગારેલા પોતાના \VUgras{સફેદ પોશાકમાં} મોહક લાગે છે.

%%K t04-c1-11-1 p
11. --- તેના \VUgras{સસરા} \textsuperscript{(20)} આટલી સૌમ્ય
\VUgras{પુત્રવધૂ} મળવાથી ઘણા સુખી છે, અને યોહાનેસનાં \VUgras{સાસુ}
\textsuperscript{(21)} પોતાની \VUgras{દીકરીને} આટલી સુંદર જોઈને સ્મિત
કરે છે.

%%K t04-c1-12-1 p
12. --- મેજને છેડે બાકીનાં \VUgras{આમંત્રિતો} \textsuperscript{(22)}
બેઠાં છે, \VUgras{સગાં, મિત્રો, કાકાના દીકરા, કાકાની દીકરીઓ.}
બગલમાં નેપકિન લઈને \VUgras{મુખ્ય પીરસનારો} \textsuperscript{(23)} તેમને
ખંતથી પીરસે છે.

%%K t04-tit-4 sub
\VUpk{10.2pt}{40}{મિત્રો.}
\VUsaut{3.0mm}

%%K t04-c1-13-1 p
13. --- મેજની જમણી બાજુએ, આ રહી \VUgras{કુમારિકા}
\textsuperscript{(24)}, જુવાન વહુની \VUgras{સખી}, પછી તેનો ભાઈ, જે
તેનો \VUgras{માનનીય સાથી} \textsuperscript{(25)} છે. તે પોતાની
\VUgras{માનનીય સખી} \textsuperscript{(26)} સાથે વાતો કરે છે, જે વરની
બહેન છે અને જે આ લગ્નથી માર્થાની \VUgras{નણંદ} બને છે. તે મોહક જુવાન
છે. તેની પાસે સુંદર \VUgras{ઘરેણાં,} \VUgras{મોતીનો હાર} અને સોનાનું
\VUgras{કંગન} છે.

%%K t04-c1-14-1 p
14. --- તેમની પાસે, જે સજ્જન ઊભા થઈને પોતાનો પ્યાલો ઊંચકે છે તે
કુટુંબના \VUgras{મિત્ર} \textsuperscript{(27)} છે. તે \VUgras{વરના
સાક્ષીઓમાંના} એક છે. તે \VUgras{માનપાન} કરે છે, વરકન્યાની
\VUgras{તંદુરસ્તી} માટે પીએ છે અને તેમને લાંબા સુખની શુભેચ્છા આપે છે.
તેમણે વિધિનો પોશાક પહેર્યો છે, પૂંછડીવાળો કોટ અને \VUgras{ચકચકતાં
ચંપલ} \textsuperscript{(28)}. તે \VUgras{દાદીમાના}
\textsuperscript{(29)} સાથી છે, જે \VUgras{વિધવા} છે.

%%K t04-c1-15-1 p
15. --- \VUgras{પૂંછડીવાળા કોટમાં} \textsuperscript{(31)} જે જુવાન છે,
પાતળી કાળી મૂછવાળો, તે યોહાનેસનો \VUgras{બનેવી} \textsuperscript{(30)}
છે. તે નામાંકિત ઇજનેર છે. તે ઘણી સુઘડ \VUgras{જુવાન સ્ત્રીના}
\textsuperscript{(33)} પડોશી છે, જે માર્થાની \VUgras{મોટી બહેન} છે અને
તે નાનકડી સુંદર છોકરીની મા છે, જે પોતાના કાકાના દીકરાનો હાથ પકડીને
ફૂલ આપવા અને તેને \VUgras{શુભ સંધ્યા} \VUgras{કહેવા} આવે છે. આ
સ્ત્રીને ઘણાં સુંદર \VUgras{કાનનાં ઘરેણાં} અને સુઘડ \VUgras{માથાનો
શણગાર} છે.

%%K t04-c1-16-1 p
16. --- કાકાનો નાનો દીકરો, પોતાના \VUgras{ઝભલા}
\textsuperscript{(36)} અને પોતાની ફૂલેલી \VUgras{ચડ્ડીને}
\textsuperscript{(37)} લીધે એટલો વિચિત્ર, ઘણો દયામણો છે. તે
\VUgras{અનાથ} \textsuperscript{(35)} છે. ઘણી નાની ઉંમરે તેણે પોતાના
પિતા અને પોતાની માતા ગુમાવ્યાં. તેના કાકા તેના \VUgras{વાલી}
\textsuperscript{(38)} છે, અને તે બિચારા બાળકને ઉછેરવા માટે કુંવારા
રહ્યા. તે ભલા માણસ છે. તેમને થોડું ટાલ છે અને તે ઘણું ઓછું જુએ છે ; તે
\VUgras{ચશ્માં} વાપરે છે.

%%K t04-tit-5 sub
\VUsaut{2.4mm}
\VUpk{10.2pt}{40}{મીઠાઈ.}
\VUsaut{3.0mm}

%%K t04-c1-17-1 p
17. --- ભોજન કરનારાંની પાછળ, \VUgras{મેજપોશથી} ઢાંકેલા મેજ પર,
\VUgras{મીઠાઈ} \textsuperscript{(39)} તૈયાર છે. મોટા પ્યાલામાં, આ રહ્યાં
ફળ, \VUgras{આલૂ} \textsuperscript{(40)}, \VUgras{નાસપતી}
\textsuperscript{(41)}, \VUgras{સફરજન} \textsuperscript{(42)},
\VUgras{જરદાળુ} \textsuperscript{(43)} અને \VUgras{દ્રાક્ષ}
\textsuperscript{(44)}. પાસે, \VUgras{અનનાસ} \textsuperscript{(45)},
સુંદર \VUgras{કેક} \textsuperscript{(46)}, \VUgras{દારૂની બાટલીઓ}
\textsuperscript{(48)} અને \VUgras{શેમ્પેન} \textsuperscript{(49)} અને
સ્વચ્છ \VUgras{થાળીઓની} \textsuperscript{(47)} થપ્પીઓ પણ.

%%K t04-c1-18-1 p
18. --- \VUgras{ભંડારી} \textsuperscript{(50)} \VUgras{બાટલી}
\textsuperscript{(52)} ખોલે છે, જે જૂના દ્રાક્ષાસવની છે, પોતાના
\VUgras{બૂચકાઢણાથી} \textsuperscript{(51)}. \VUgras{બૂચ}
\textsuperscript{(53)} ઘણો કઠણ ખેંચાય એવો લાગે છે. આ ભંડારીની બગલમાં
\VUgras{નેપકિન} \textsuperscript{(54)} છે.

%%K t04-c1-19-1 p
19. --- રાતના ભોજન પછી, પુરુષો ધૂમ્રપાનખંડમાં જશે, અને સ્ત્રીઓ
નૃત્ય માટે તૈયાર થવા જશે.

%%K t04-tit-6 sub
\VUsaut{5.0mm}
\VUcentre{11.4pt}{0}{\textit{બીજું દૃશ્ય.}}
\VUsaut{1.6mm}
\VUpk{11.4pt}{40}{દાદાનો જન્મદિવસ.}
\VUsaut{2.6mm}

%%K t04-tit-7 sub
\VUpk{10.2pt}{40}{મુલાકાત.}
\VUsaut{3.0mm}

%%K t04-c2-01-1 p
1. --- દસ વર્ષ વીતી ગયાં, યોહાનેસનાં માતાપિતા વૃદ્ધ થયાં છે અને
પોતાનાં ત્રણ પૌત્રોને ઘણો પ્રેમ કરે છે, બે \VUgras{પૌત્રો}
\textsuperscript{(2)} અને એક \VUgras{પૌત્રી} \textsuperscript{(3)}. આજે
\VUgras{દાદાનો જન્મદિવસ} હોવાથી, તેમનાં સંતાનો તેમને શુભેચ્છા આપવા
આવે છે.

%%K t04-c2-02-1 p
2. --- નાનકડો પેત્રુસ હજી ઘણો નાનો છે, તે માત્ર ત્રણ વર્ષનો છે. તે
\VUgras{બિલાડીને} \textsuperscript{(1)} નજીક આવતી જુએ છે. તે તેની સુંદર
રેશમી \VUgras{રુવાંટી} અને તેની લાંબી \VUgras{પૂંછડી} પંપાળવા માગે છે,
તે વિચારતો નથી કે તેના સુંદર \VUgras{પંજા} નીચે અણીદાર દુષ્ટ નહોર
છુપાયેલા છે, જે કદાચ તેને ઉઝરડશે.

%%K t04-c2-03-1 p
3. --- તેની બહેન ગાબ્રિએલા, જે તેને હાથ આપે છે, ઘણી વધારે ગંભીર છે,
અને માર્સેલુસ પોતાના મોટા \VUgras{ડગલા} \textsuperscript{(4)} અને
ચામડાના \VUgras{પટ્ટા} \textsuperscript{(5)} સાથે થોડો મોટો થયેલો લાગે
છે, જોકે તેની ટૂંકી ચડ્ડી તેનાં \VUgras{મોજાં} \textsuperscript{(6)}
દેખાડે છે.

%%K t04-c2-04-1 p
4. --- તેમના પિતાએ પોતાનો શિયાળુ જાડો \VUgras{ઓવરકોટ}
\textsuperscript{(7)} પહેર્યો છે, જેને \VUgras{રુવાંટીનો કોલર}
\textsuperscript{(8)} છે, અને પોતાના \VUgras{ખિસ્સામાં}
\textsuperscript{(9)} તેમની પાસે રૂમાલ છે. તેમની પત્નીએ પોતાની નમદાની
ટોપી પહેરી છે, જેને \VUgras{પીંછાં} \textsuperscript{(10)} શણગારે છે,
પોતાનું \VUgras{જાકીટ} \textsuperscript{(11)}, પોતાની \VUgras{રુવાંટીની
પટ્ટી} \textsuperscript{(12)} અને પોતાનું \VUgras{હાથગરમણું}
\textsuperscript{(13)}.

%%K t04-c2-05-1 p
5. --- ઘણી ઠંડી છે, કેમ કે હેમંત ચાલે છે. \VUgras{બરફ}
\textsuperscript{(19)} આખો દિવસ પડ્યો છે અને ઘરોનાં છાપરાં તદ્દન સફેદ
છે. \VUgras{રાત} સાથે ઠંડી વધારે તીવ્ર બને છે. આકાશમાં
\VUgras{ચંદ્ર} \textsuperscript{(17)} અને \VUgras{તારા}
\textsuperscript{(18)} ચમકે છે અને \VUgras{અંધકારને} ભેદે છે.

%%K t04-tit-8 sub
\VUsaut{4.2mm}
\VUpk{10.2pt}{40}{માંદગી.}
\VUsaut{3.0mm}

%%K t04-c2-06-1 p
6. --- બિચારો \VUgras{આંધળો} \textsuperscript{(20)} ઘણો કમનસીબ છે, જે
\VUgras{દવાની દુકાન} \textsuperscript{(16)} આગળનું ચોગાન ઓળંગે છે અને
જેને પોતાનો \VUgras{કૂતરો} \textsuperscript{(21)} દોરે છે. યુસ્તિનિયા,
\VUgras{નોકરાણી} \textsuperscript{(14)}, રસોડામાંથી \VUgras{વાટકો}
\textsuperscript{(15)} લાવે છે, જે ઉદારતાથી ખાંડ નાખેલા ઘણા ગરમ
\VUgras{ઉકાળાનો} છે.

%%K t04-c2-07-1 p
7. --- \VUgras{દાદીમા} \textsuperscript{(23)}, જે મેજ પરથી પોતાનાં
\VUgras{દાંડીવાળાં ચશ્માં} \textsuperscript{(26)} શોધવા માગે છે, જેથી
તે \VUgras{ચિઠ્ઠી} \textsuperscript{(27)} વાંચી શકે, જે
\VUgras{ડૉક્ટરે} \textsuperscript{(25)} હમણાં લખી છે, પોતાની
\VUgras{લાકડી} \textsuperscript{(24)} પર ટેકવીને પોતાની મોટી
ખુરશીમાંથી ઊભાં થાય છે.

%%K t04-c2-08-1 p
8. --- ઘણી વાર તેમને \VUgras{નાની બીમારીઓ, ચક્કર, શરદી, દાંતના દુખાવા}
થાય છે, તેમના પગ નબળા છે અને તેમની આંખો હવે બહુ સારી નથી. તેમ છતાં,
તે ખુશીથી પોતાના જૂના \VUgras{ડૉક્ટરની} મજાક કરે છે, જે તેમને હંમેશાં
\VUgras{દવા} \textsuperscript{(28)} લખી આપવા માગે છે. આજે તે ઘણાં ખુશ છે
કેમ કે તેમની આસપાસ તેમનું જુવાન કુટુંબ છે.

%%K t04-c2-09-1 p
9. --- સારી રીતે બંધ કરેલા ઓરડામાં આરામ લાગે છે. આરસની
\VUgras{સગડીમાં} \textsuperscript{(29)} \VUgras{સુંદર આગ}
\textsuperscript{(30)} બળે છે, અને હમણાં જ સળગાવેલા \VUgras{દીવાના}
\textsuperscript{(31)} મીઠા \VUgras{પ્રકાશમાં} માણસ સુખી લાગે છે.

%%K t04-tit-9 sub
\VUsaut{4.2mm}
\VUpk{10.2pt}{40}{દાદા.}
\VUsaut{3.0mm}

%%K t04-c2-10-1 p
10. --- \VUgras{દાદા} \textsuperscript{(33)} પોતે પણ ભૂલી ગયા છે કે તે
માંદા હતા, કે તેમનો \VUgras{વા} તેમને પોતાની \VUgras{મોટી ખુરશીમાં}
\textsuperscript{(34)} જકડી રાખે છે અને કે ગઈ કાલે પણ તેમને
\VUgras{તાવ અને બેભાનપણું} હતું. તેમને એ પણ યાદ નથી કે ગયા શિયાળામાં
તેમને \VUgras{ફેફસાંનો સોજો} થયો હતો અને કે કર્કશ અને પીડાદાયક
\VUgras{ઉધરસે} તેમને ઘણું દુખ આપ્યું હતું.

%%K t04-c2-10-2 p
તેમ છતાં ઘણી મુશ્કેલીથી તે સાજા થયા. હવે તે થોડા સારા છે. પણ તેમનો પગ,
જે તે \VUgras{ઓશીકા} \textsuperscript{(38)} પર ટેકવે છે, જે નાના
\VUgras{ટેભા} \textsuperscript{(39)} પર મૂક્યું છે, તેમને પણ દુખે છે.

%%K t04-c2-11-1 p
11. --- જ્યારે તેમને કાળજીથી પોતાના હૂંફાળા \VUgras{ઝભ્ભામાં}
\textsuperscript{(35)} વીંટવામાં આવે છે, જેને છીપનાં જાડાં
\VUgras{બટન} \textsuperscript{(36)} છે, અને જ્યારે તેમની પાસે તેમને
છાપું વાંચી સંભળાવવા માટે નાનકડી \VUgras{અનાથ છોકરી}
\textsuperscript{(40)} હોય છે, જેણે સફેદ \VUgras{ટોપી}, ભૂરો
\VUgras{ઝભ્ભો} \textsuperscript{(42)} અને \VUgras{ખભાવસ્ત્ર}
\textsuperscript{(41)} પહેર્યાં છે, ત્યારે તે ફરિયાદ કરતા નથી.

%%K t04-c2-12-1 p
12. --- દર વર્ષે, જ્યારે \VUgras{કૅલેન્ડર} \textsuperscript{(43)} ફરી
પહેલી \VUgras{ડિસેમ્બરની} \VUgras{તારીખ} બતાવે છે, ત્યારે તે અધીરાઈથી
પોતાનાં \VUgras{પૌત્રોની} રાહ જુએ છે, કેમ કે તે જાણે છે કે તેઓ આ
મહત્ત્વની \VUgras{તારીખ} ભૂલશે નહીં.

%%K t04-c2-13-1 p
13. --- તે લાગણીથી ઘણા જૂના સમયને યાદ કરે છે, જ્યારે તે પોતે
યશસ્વી શાહી \VUgras{સેનાપતિને} શુભેચ્છા આપવા ગયા હતા, જેમનું
\VUgras{ચિત્ર} \textsuperscript{(44)} ભીંતે લટકે છે અને જે તેમનાં
સંતાનોના \VUgras{પરદાદા} છે.

\VUsaut{6.0mm}
\VUfilet{22mm}
